\section{Antes de empezar}
\subsection*{Cómo entras}
\begin{enumerate}
  \item \emph{Titulación} → \emph{Posgrado: validar trámites} (bandeja).
  \item Los datos que salen en los oficios (quién firma, contacto) se editan en \emph{Posgrado} → \emph{Configuración}.
\end{enumerate}
\begin{tencuenta}
\begin{itemize}
  \item El cotejo solo se registra con todos los requisitos validados o recibidos.
  \item Programar con firmas faltantes, respuestas «No» o vigencias que no alcanzan exige escribir una justificación, que queda en la bitácora.
  \item Un resultado «Aplazado» reabre el segundo momento para reprogramar.
\end{itemize}
\end{tencuenta}
\begin{algosalemal}
\begin{itemize}
  \item Si el egresado necesita corregir sus datos antes del cotejo, anula el documento firmado con una justificación.
  \item Si un sinodal contestó «No» en el visto bueno, puedes devolver la propuesta al egresado; se cancelan las firmas.
  \item Si faltan acuses, puedes cerrar las invitaciones con una justificación.
\end{itemize}
\end{algosalemal}
